\section{Conserved occupations and explicit critical bounds}
\label{sec:r13-occupations}

The local critical nonlinearity admits an explicit connection between
scalar sharp interpolation and the entire conserved occupation spectrum.
The auxiliary Gaussian variables below occur only in the proof; the
initial data and the evolution remain deterministic.

Let $\mathcal H$ be a separable complex Hilbert space and suppose the
scalar inequality is
\begin{equation}\label{eq:r13-scalarGN}
 \int |f|^{2(1+r)}\le (1+r)M_c^{-r}M(f)^r A(f),
 \qquad r=m/N>0.
\end{equation}
The positive scalar Fourier multiplier defining $A$ may be anisotropic.
For nonzero $U\in H^{m/2}(\mathbb R^N;\mathcal H)$ set
\[
 \mathcal F(U)=\int\|U\|_{\mathcal H}^{2(1+r)},\qquad
 \Gamma=\int|U(x)\rangle\langle U(x)|\,dx,\qquad
 \Gamma_A=\int a(\xi)|\widehat U(\xi)\rangle
                              \langle\widehat U(\xi)|\,d\xi.
\]
Write $M=\operatorname{Tr}\Gamma$, $A=\operatorname{Tr}\Gamma_A$,
$\mu=\|\Gamma\|_{\rm op}$, and list the positive eigenvalues of
$\Gamma$ as $\lambda_1=\mu\ge\lambda_2\ge\cdots$.
Euler's gamma function is denoted by $\mathrm G$ in this section, to
distinguish it from the correlation operator.

\begin{theorem}[Explicit bound from the occupation spectrum]
\label{thm:r13-Gaussian}
Let $\xi_j$ be independent exponential random variables of mean one.
Define
\begin{equation}\label{eq:r13-Bdef}
 \mathcal B_r(\Gamma)=\frac{1}{\mathrm G(r+2)}
   \mathbb E\left[\xi_1\left(\sum_j\lambda_j\xi_j\right)^r\right].
\end{equation}
Then
\begin{equation}\label{eq:r13-spectralGN}
 \mathcal F(U)\le(1+r)M_c^{-r}\mathcal B_r(\Gamma)A(U).
\end{equation}
The coefficient depends only on the spectrum, is homogeneous of degree
$r$, and satisfies $\mathcal B_r(\Gamma)\le M^r$. Equality in this last
inequality holds precisely when $\Gamma$ has rank one.
\end{theorem}
\begin{proof}
First suppose the internal rank is finite and diagonalize $\Gamma$.
For independent standard circular complex Gaussian variables $z_j$, put
$f_z=\sum_jz_ju_j$. Pointwise, $f_z$ has Gaussian variance
$\rho_U=\|U\|^2$, so
\[
 \mathbb E\mathcal F(f_z)=\mathrm G(r+2)\mathcal F(U).
\]
In the chosen basis $M(f_z)=S:=\sum_j\lambda_j|z_j|^2$.
Apply \eqref{eq:r13-scalarGN} and average. Averaging over the phases
of the $z_j$ removes off-diagonal entries of the kinetic Gram matrix,
giving
\begin{equation}\label{eq:r13-weightedGaussian}
 \mathrm G(r+2)\mathcal F(U)
 \le(1+r)M_c^{-r}\sum_j A(u_j)b_j,
 \qquad b_j=\mathbb E[\xi_j S^r].
\end{equation}
If $\lambda_i\ge\lambda_j$, exchange $\xi_i$ and $\xi_j$ in one
copy of this expectation. With $S'$ denoting the exchanged sum,
\[
 b_i-b_j=\tfrac12\mathbb E[(\xi_i-\xi_j)(S^r-(S')^r)]\ge0,
\]
because $S-S'=(\lambda_i-\lambda_j)(\xi_i-\xi_j)$.
Thus $b_1$ is maximal, proving \eqref{eq:r13-spectralGN}.
H\"older's inequality and convex Jensen at exponent $r+1$ give
\[
 \mathbb E[\xi_1S^r]
 \le(\mathbb E\xi_1^{r+1})^{1/(r+1)}
       (\mathbb ES^{r+1})^{r/(r+1)}
 \le\mathrm G(r+2)M^r.
\]
If a second eigenvalue is positive, $S$ is not proportional to $\xi_1$
almost surely, and H\"older is strict. In rank one the moment is
exactly $M^r\mathrm G(r+2)$.

For infinite rank, truncate in the eigenbasis of $\Gamma$.
The truncated fields converge in $H^{m/2}$ and in the interaction
space. Their moments increase to \eqref{eq:r13-Bdef} by monotone
convergence and remain bounded by $M^r$. This proves the full result.
In particular no identity-covariance Gaussian vector in infinite-dimensional
$\mathcal H$ is assumed to exist.
\end{proof}

\subsection{Explicit quantitative polarization}
For $X=\Gamma/M$ and $Y=\Gamma_A/A$, both positive of trace one, the
normalized distance to a fixed internal direction is exactly
\begin{equation}\label{eq:r13-pdistance}
 \mathfrak p(U):=\inf_{\|\zeta\|=1,q}
 \left\{\frac{M(U-\zeta q)}M+\frac{A(U-\zeta q)}A\right\}
 =2-\lambda_{\max}(X+Y).
\end{equation}
Indeed the pointwise scalar projection minimizes the combined quadratic
norm for each $\zeta$.

\begin{theorem}[Mass and kinetic alignment, $r\ge1$]
\label{thm:r13-explicitpolarization}
For complex Hilbert-valued fields,
\begin{align}
 \frac{\mathcal F(U)}{(1+r)(M/M_c)^r A}
 &\le\frac{1+r\operatorname{Tr}(XY)}{r+1},
 \label{eq:r13-alignment}\\
 \mathfrak p(U)&\le\frac{2(r+1)}r\,\mathfrak d(U),
 \qquad
 \mathfrak d(U)=1-
  \frac{\mathcal F(U)}{(1+r)(M/M_c)^r A}.
 \label{eq:r13-explicitpol}
\end{align}
Consequently
\begin{equation}\label{eq:r13-simpleoccupation}
 \mathcal F(U)\le(1+r)(M/M_c)^r
 \left[1-\frac r{r+1}\left(1-\frac\mu M\right)\right]A.
\end{equation}
For real Hilbert-valued fields with real-even dispersion, the respective factors in
\eqref{eq:r13-alignment} and \eqref{eq:r13-explicitpol} are
\[
 \frac{1+2r\operatorname{Tr}(XY)}{2r+1},
 \qquad \frac{2r+1}{r}.
\]
\end{theorem}
\begin{proof}
For $r\ge1$, Jensen gives
$S^r\le M^{r-1}\sum_i\lambda_i\xi_i^r$. Hence
\[
 \frac{b_j}{\mathrm G(r+2)}
 \le M^{r-1}\left[\lambda_j+\frac{M-\lambda_j}{r+1}\right].
\]
Substitution into \eqref{eq:r13-weightedGaussian} proves
\eqref{eq:r13-alignment}. Since
\[
 \lambda_{\max}(X+Y)
 \ge\frac{\operatorname{Tr}((X+Y)^2)}{\operatorname{Tr}(X+Y)}
 \ge2\operatorname{Tr}(XY),
\]
where the second step is $\operatorname{Tr}((X-Y)^2)\ge0$,
\eqref{eq:r13-pdistance} proves \eqref{eq:r13-explicitpol}.
Also $\operatorname{Tr}(XY)\le\|X\|_{\rm op}=\mu/M$, proving
\eqref{eq:r13-simpleoccupation}.
For real Gaussian projections use $\xi_j=g_j^2$ with standard real
$g_j$. Their moments satisfy
\[
 a_q=\mathbb E\xi_j^q
 =\frac{2^q\mathrm G(q+1/2)}{\sqrt\pi},
 \qquad a_{r+1}=(2r+1)a_r.
\]
The same calculation gives the stated real constants. Finite-rank
approximation proves all assertions in infinite dimension.
\end{proof}

A scalar complex field is a real two-component field
$(\operatorname{Re}u,\operatorname{Im}u)$. Therefore, for $r\ge1$,
there is a constant phase $\theta$ and a real scalar $q$ such that
\begin{equation}\label{eq:r13-explicitphase}
 \frac{M(u-e^{i\theta}q)}{M(u)}+
 \frac{A(u-e^{i\theta}q)}{A(u)}
 \le\frac{2r+1}{r}\mathfrak d(u).
\end{equation}
This supplies the phase-smallness input for near-critical quartic
motion in $N\le4$ directly from sharp scalar interpolation.
The real covariance used for this instantaneous estimate is generally
not conserved. The complex internal covariance in the following
dynamical result is conserved.

\subsection{Consequences for the actual evolution}
Consider the same common scalar system
\[
 i\partial_tU=(a+b)(D)U-\|U\|^{2r}U,\qquad b\ge0,
 \qquad E=\frac{A+B}{2}-\frac{\mathcal F}{2(1+r)}.
\]
Assume its conservative energy-space local flow and continuation
property. All pairwise component inner products, and therefore
$\Gamma$, are conserved. To see this, use the common scalar
self-adjoint Hamiltonian $H_U=(a+b)(D)-\|U\|^{2r}$. For smooth
approximations the equations give
\[
 \frac{d}{dt}\int u_j\overline{u_k}
 =-i\int (H_Uu_j)\overline{u_k}
   +i\int u_j\overline{H_Uu_k}=0.
\]
The quadratic-form or strong-approximation interpretation passes this
identity to the energy-space flow. In infinite internal dimension,
all matrix elements are conserved and determine the positive trace-class
operator $\Gamma$; continuity also follows from
$\|\Gamma(U)-\Gamma(V)\|_1
\le(\|U\|_2+\|V\|_2)\|U-V\|_2$.

\begin{theorem}[Global bounded flow from conserved occupations]
\label{thm:r13-spectrumflow}
If $\mathcal B_r(\Gamma(U_0))<M_c^r$, the solution is global in both
time directions, and
\[
 \sup_t A(U(t))\le
 \frac{2E(U_0)}{1-\mathcal B_r(\Gamma(U_0))/M_c^r},
 \qquad \sup_t B(U(t))\le2E(U_0).
\]
In particular any fixed normalized occupation spectrum of rank greater
than one has an explicitly specified global mass interval extending
strictly above $M_c$. At mass $M_c$, every orbit with unbounded energy
norm lies in one fixed internal line for its entire maximal evolution.
\end{theorem}
\begin{proof}
Theorem~\ref{thm:r13-Gaussian} and conservation of $\Gamma$ yield
\[
 E(U(t))\ge\frac12
 \left[1-\frac{\mathcal B_r(\Gamma(U_0))}{M_c^r}\right]A(U(t))
 +\frac12 B(U(t)).
\]
Conservation and continuation prove the bounds and global existence.
Homogeneity and strictness of $\mathcal B_r(\Gamma)<M^r$ give the
above-critical interval. At $M=M_c$, unboundedness forces rank one.
The fixed range of the conserved $\Gamma$ then contains $U(t,x)$
almost everywhere at every time.
\end{proof}

Proposition~\ref{prop:r13g-local-flow} verifies these flow assumptions
for $m=2k\ge4$ and $N=1,2,3$, also for countably many components.
In particular it applies to quartic dispersion in these dimensions.
No randomness of the initial field,
radial symmetry, or finite spatial moment is required.

The spectral moment is directly computable. For integer $r$, with
$h_r$ the complete homogeneous symmetric polynomial,
\[
 \mathcal B_r(\Gamma)
 =\frac{h_r(\lambda_1,\lambda_1,\lambda_2,\ldots)}{r+1}.
\]
In particular,
\begin{equation}\label{eq:r13-B2}
 \mathcal B_2(\Gamma)=
 \frac{M^2+\operatorname{Tr}\Gamma^2+2M\mu+2\mu^2}{6}.
\end{equation}
For $k$ equal occupations $\lambda_j=M/k$,
\[
 \frac{\mathcal B_r(\Gamma)}{M^r}
 =\frac{\mathrm G(k+r+1)}
 {k^r\mathrm G(k+1)\mathrm G(r+2)}.
\]
Thus in the two-dimensional quartic system, two equal orthogonal
occupations give $\mathcal B_2=M^2/2$: every such $H^2$ datum with
$M<\sqrt2M_c$ has globally bounded evolution.
Formula~\eqref{eq:r13-B2} retains information beyond $M$ and $\mu$.
For $r\ge1$, at fixed $M$ and $\lambda_1$, averaging any two of the
remaining occupations cannot increase $\mathcal B_r$, by convexity
and permutation symmetry of its defining expectation.

These estimates retain the sharp scalar endpoint, but no assertion
of optimality at each higher-rank spectrum is made. Evaluating that
optimal coefficient is a sharp orthonormal kinetic problem. The
Gaussian argument above applies to local powers; Hartree interactions
produce an additional exchange term under Gaussian averaging, so no
Hartree extension is asserted here.
